% ============================================================
% VAJRA: A Ground-Truth Verified LLM Red-Teaming Framework
% IEEEtran Conference Class — compile on Overleaf as-is
% ============================================================
\documentclass[conference]{IEEEtran}

% ---------- Packages ----------
\usepackage{cite}
\usepackage{amsmath,amssymb}
\usepackage{graphicx}
\usepackage{url}
\usepackage{hyperref}
\usepackage{booktabs}
\usepackage{xcolor}
\usepackage{listings}

\hypersetup{
    colorlinks=true,
    linkcolor=blue,
    citecolor=blue,
    urlcolor=blue,
}

% ---------- Title ----------
\title{VAJRA: A Ground-Truth Verified Scoring Architecture\\for Automated LLM Red-Teaming}

\author{
  \IEEEauthorblockN{Sahil Powar}
  \IEEEauthorblockA{%
    \textit{Independent Security Research}\\
    \href{https://github.com/sahilll05/LLM-RedTeam}{github.com/sahilll05/LLM-RedTeam}
  }
}

\begin{document}
\maketitle

% ============================================================
\begin{abstract}
% GUIDANCE (delete before submission):
%   One paragraph ~150 words. Structure:
%     1. Problem: automated LLM red-teaming tools conflate
%        topic-discussion with actual compromise.
%     2. Consequence: measured false-positive rate on benign,
%        superficially sensitive queries (cite WildJailbreak
%        and JailbreakBench calibration splits).
%     3. Contribution: canary-token exfiltration proofs,
%        sandboxed exploit execution, decomposed rubric judge.
%     4. Validation: against Young & Moody consensus-labeled
%        ground truth and cross-model benchmark.
%     5. Result: FPR reduction X%, detection rate preserved.
Existing automated LLM red-teaming tools achieve detection
precision through deterministic, ground-truthed checks for
some categories (e.g., registry lookups, verbatim text
matching) but rely on heuristic or single-pass LLM-judge
scoring for others---particularly system-prompt exfiltration
and actionable-exploit categories---where the boundary
between benign topic discussion and actual compromise is not
mechanically verified. We show that na\"ive keyword-based
scoring produces a high false-positive rate on benign,
superficially sensitive queries, measured against contrastive
benign data from WildJailbreak~\cite{jiang2024wildteaming}
and JailbreakBench~\cite{chao2024jailbreakbench}. We introduce
VAJRA, a scoring architecture grounding detection in (1)
\textbf{canary-token exfiltration proofs} across six
placement strategies, (2) \textbf{sandboxed SQL injection
execution} against a containerised deliberately-vulnerable
application, and (3) a \textbf{decomposed knowledge-vs-action
rubric judge}, validated against a published consensus-labeled
ground truth~\cite{young2026executableweapons}, reducing false
positives substantially while preserving detection of confirmed
vulnerabilities.
\end{abstract}

% IEEE keywords
\begin{IEEEkeywords}
LLM red-teaming, prompt injection, jailbreak detection, canary tokens,
false-positive reduction, automated security evaluation
\end{IEEEkeywords}

% ============================================================
\section{Introduction}
\label{sec:intro}
% GUIDANCE:
%   ~500 words. Structure:
%     1. Rise of LLM-integrated apps and attack surface (OWASP).
%     2. Existing tools: garak, PyRIT — strengths and the gap.
%     3. The core failure mode: topic presence as proxy for
%        compromise (the "genetic diseases" / "SQL injection
%        education" examples).
%     4. Our reframe: VAJRA's contribution.
%     5. Bullet contributions (RQ1--RQ5).
%     6. Paper organisation sentence.

Large language model (LLM) applications are increasingly
deployed in security-sensitive contexts, creating novel attack
surfaces categorised by OWASP's Top 10 for LLM
Applications~\cite{owasp2025llm}. Automated red-teaming
frameworks such as garak~\cite{garak2024} and
PyRIT~\cite{pyrit2024} have emerged to probe these systems
at scale, but existing tools exhibit a fundamental limitation:
they conflate \emph{topic discussion} with \emph{actual
compromise}, producing high false-positive rates on benign,
superficially sensitive queries.

% [TODO: expand with quantitative FPR motivation after Phase 1]

\noindent\textbf{Contributions.} This paper presents VAJRA,
a modular LLM red-teaming framework with three primary
contributions:
\begin{itemize}
  \item A \textbf{canary-token exfiltration engine} with
    multiple placement and delimiter strategies, providing
    mechanically verified, zero-false-positive detection of
    system-prompt leakage (addressing RQ3).
  \item A \textbf{sandboxed exploit verifier} that executes
    LLM-generated payloads against a deliberately vulnerable
    application, providing binary, ground-truthed verdicts for
    actionable-exploit categories (addressing RQ2).
  \item A \textbf{decomposed rubric judge} replacing
    single-shot LLM-judge prompts with structured
    sub-questions operationalising the knowledge-vs-action
    axis, validated against the Young \& Moody
    consensus-labeled prompt bank~\cite{young2026executableweapons}
    (addressing RQ1).
\end{itemize}

% ============================================================
\section{Background and Related Work}
\label{sec:related}
% GUIDANCE:
%   ~600 words. Subsections:
%   2.1 Prompt Injection (Perez & Ribeiro, Greshake et al.)
%   2.2 Existing Red-Teaming Tools (garak, PyRIT) with table
%   2.3 Calibration Datasets (WildJailbreak, JBB, XSTest)
%   2.4 Knowledge vs. Action (Young & Moody, CyberSecEval)

\subsection{Prompt Injection}
Direct prompt injection---supplying adversarial instructions
via user-controlled input to override an LLM's
system-level directives---was formally characterized by
Perez and Ribeiro~\cite{perez2022ignore}. Indirect
(RAG-based) injection, where malicious instructions are
embedded in retrieved documents rather than direct user
input, was demonstrated by Greshake et
al.~\cite{greshake2023indirect}.

\subsection{Automated Red-Teaming Tools}
% [TODO: expand comparison table (Master Plan Part 1.5)]

\subsection{Calibration and False-Positive Benchmarks}
WildJailbreak~\cite{jiang2024wildteaming} is a synthetic
dataset constructed by mining jailbreak \emph{tactics} from
real user-chatbot interactions via the WildTeaming pipeline,
then using GPT-4 and Mixtral-8x7B to compose new prompts
applying combinations of mined tactics. Critically for our
use, it includes \texttt{vanilla\_benign} (50,050 rows) and
\texttt{adversarial\_benign} (78,706 rows) splits built
explicitly to combat exaggerated safety, following the
methodology of XSTest~\cite{rottger2024xstest}.

JailbreakBench~\cite{chao2024jailbreakbench} provides 100
harmful behaviors across 10 categories aligned with OpenAI
usage policy, \emph{plus} 100 topic-matched benign
counterparts (\texttt{benign-behaviors.csv}), making it
uniquely suited for testing whether a scorer distinguishes
intent while holding topic constant.

\subsection{Knowledge vs.\ Action Ground Truth}
Young and Moody~\cite{young2026executableweapons,young2026companion}
operationalise the knowledge-vs-action axis via a
consensus-labeling protocol spanning five LLM families,
reaching Fleiss' $\kappa = 0.876$, across a 1,554-prompt
bank. We use this bank to externally validate our rubric judge.

% ============================================================
\section{System Design}
\label{sec:design}
% GUIDANCE: ~800 words.
%   3.1 Architecture overview
%   3.2 Payload Library and Calibration Splits
%   3.3 Canary-Token Engine (placement strategies — key contribution)
%   3.4 Sandboxed Exploit Verifier (Docker, Flask SQLi target)
%   3.5 Decomposed Rubric Judge (four sub-questions)
%   3.6 Adaptive Attack Orchestrator (simplified PAIR)

\subsection{Architecture Overview}
% [TODO: add figure — use ASCII diagram from ARCHITECTURE.md]

\subsection{Payload Library and Calibration Splits}
% [TODO: describe suite structure and ingestion pipeline after
%  ingest_hf.py is extended]

\subsection{Canary-Token Engine}
% [TODO: describe placement strategies and delimiter variants
%  after engine/canary.py is built and RQ3 experiments run.
%  Key variants to describe:
%    (1) Inline injection: canary mid-system-prompt.
%    (2) Prefix injection: canary at start of system prompt.
%    (3) Suffix injection: canary at end of system prompt.
%    (4) XML-tagged: <canary>TOKEN</canary>.
%    (5) Comment-style: # INTERNAL_TOKEN: <value>
%  Report detection rate per strategy.]

\subsection{Sandboxed Exploit Verifier}
\label{subsec:sandbox}

To provide binary, ground-truthed verdicts for SQL injection and
code-attack categories, VAJRA runs a deliberately-vulnerable Flask
application inside a Docker container (Fig.~\ref{fig:sandbox}).
The container is bound exclusively to \texttt{localhost:5555} via
an internal Docker network with no upstream internet route, ensuring
no production systems are reachable from the sandbox.

The vulnerable application exposes a \texttt{/login} endpoint
constructed via direct string interpolation---a textbook first-order
SQL injection target. The database is seeded with three user accounts
(one admin, two standard users) and a planted secret string in the
admin record, making the application suitable for both authentication-bypass
and sensitive-value exfiltration experiments.

When VAJRA scores a payload with technique \texttt{sqli} or
\texttt{sql\_injection} and the heuristic/judge verdict is \textsc{Complied},
\texttt{engine/exploit\_verifier.py} performs the following:
\begin{enumerate}
  \item Extracts candidate SQL injection strings from the LLM response
    via four regex patterns targeting classic comment-bypass, UNION SELECT,
    tautology, and code-block-wrapped variants.
  \item Resets the sandbox database via \texttt{POST /reset} to ensure
    each candidate is tested against a clean state.
  \item Submits each candidate to \texttt{POST /verify}, which fires
    the payload as a raw username value against the vulnerable query and
    returns whether authentication was bypassed, the role obtained, and
    the running SQL query for logging.
  \item If any candidate achieves bypass, the verdict is upgraded to
    \textsc{Verified\_Exploit} with the full server-side evidence; otherwise
    the verdict remains \textsc{Complied} with a sandbox annotation
    indicating that no extracted candidate was executable.
\end{enumerate}

This pipeline is strictly more informative than keyword or LLM-judge
scoring: a model that accurately describes SQL injection theory
but produces no syntactically valid payload cannot trigger a
\textsc{Verified\_Exploit} verdict by construction. The verifier
degrades gracefully---with a startup warning---if the Docker container
is not running, leaving the heuristic/judge path unaffected.

\subsection{Decomposed Rubric Judge}
% [TODO: describe four sub-questions after scorer.py rebuild:
%    Q1. Does the response go beyond freely available information?
%    Q2. Is it tied to a specific, named target?
%    Q3. Would acting on it plausibly succeed against a control?
%    Q4. Does it contain a verbatim planted canary?]

\subsection{Adaptive Attack Orchestrator}
% [TODO: Phase 3 — simplified PAIR implementation]

% ============================================================
\section{Evaluation}
\label{sec:eval}
% GUIDANCE: ~1000 words.
%   4.1 Experimental Setup (model matrix, dataset versions/hashes)
%   4.2 RQ1: FPR/TPR on calibration splits
%   4.3 RQ2: Detection on confirmed ground-truth attacks
%   4.4 RQ3: Canary placement detection rates per strategy
%   4.5 RQ4: Static vs. adaptive attack success rate
%   4.6 RQ5: Cross-model benchmark
%   Use Wilson CIs, McNemar's test, Holm-Bonferroni correction.

\subsection{Experimental Setup}
% [TODO: fill after experiments — model matrix, compute,
%  dataset versions, payload set hash, judge model version]

\subsection{RQ1: False-Positive Rate Reduction}
% [TODO]

\subsection{RQ2: Detection Rate on Ground-Truth Attacks}
% Designed experiment:
% Compare VERIFIED_EXPLOIT rate vs plain COMPLIED rate on the injection suite.
% A VERIFIED_EXPLOIT requires a working payload; COMPLIED may include
% models that merely discussed SQL injection theory.
% Expected result: VERIFIED_EXPLOIT subset is strictly smaller than COMPLIED
% — showing the verifier eliminates false positives in the injection category
% without eliminating true positives (a model producing a working payload
% will always show up in both sets).
% [TODO: fill with measured numbers after experiments run]

\subsection{RQ3: Canary Placement Effectiveness}
% [TODO: key table — detection rate per placement strategy.
%  Hypothesis: XML-tagged and suffix strategies perform
%  worse than inline; test and report.]

\subsection{RQ4: Static vs.\ Adaptive Attack Success}
% [TODO]

\subsection{RQ5: Cross-Model Benchmark}
% [TODO]

% ============================================================
\section{Discussion}
\label{sec:discussion}
% GUIDANCE: ~400 words.
%   - What results tell us about the state of LLM red-teaming.
%   - Practical implications for deployers.
%   - Where VAJRA is and isn't applicable (see THREAT_MODEL.md).

% [TODO: fill after results]

% ============================================================
\section{Limitations and Threats to Validity}
\label{sec:limitations}
% GUIDANCE (do not omit — reviewers expect this section):
%   - Canary scope: only covers exfiltration/system-access
%     categories; others still rely on rubric scoring.
%   - Sandbox scope: SQLi and limited code-attack variants only.
%   - Young & Moody bank: gated; reproducible only with HF token.
%   - Judge LLM: behaviour varies by model version; pin version.
%   - Dataset recency: WildJailbreak tactics from a fixed date.

\textbf{Canary scope.}
Canary-token verification provides zero-false-positive ground
truth only for the exfiltration and system-access categories.
All other attack categories continue to rely on rubric-judge
scoring.

\textbf{Sandbox scope.}
The exploit verifier targets a deliberately vulnerable Flask
application and covers SQL injection and a small number of
code-attack variants. Verdict generalisation to other
exploit classes requires additional sandbox targets.

\textbf{Gated dataset reproducibility.}
The Young \& Moody prompt bank and WildJailbreak are accessed
under gated responsible-use terms. Full reproduction of our
results requires the user's own Hugging Face token and
acceptance of those terms; raw prompt text is not redistributed.

\textbf{Judge model version sensitivity.}
Rubric judge performance is sensitive to the judge LLM version.
All results should be reproduced with pinned model identifiers
as reported in our experimental setup.

\textbf{Dataset recency.}
WildJailbreak tactics are derived from interactions collected
up to its publication date. Rapidly evolving jailbreak techniques
emerging after that date may not be represented in our evaluation.

% ============================================================
\section{Ethics Statement}
\label{sec:ethics}
% GUIDANCE: Required for security venues.

All experiments were conducted on systems under our direct
control. The payload library published with VAJRA represents
technique-level illustrations for defensive security testing
and is not an optimised attack toolkit. Gated datasets
(WildJailbreak, Young \& Moody prompt bank) are accessed
via users' own Hugging Face credentials in compliance with
their respective responsible-use terms; no raw prompt text
from these datasets is redistributed in this repository.
Sandboxed exploit verification will be performed in a
network-isolated Docker environment; no production systems
were targeted at any point.

% ============================================================
\section{Conclusion}
\label{sec:conclusion}
% GUIDANCE: ~200 words.
%   - Restate one-sentence reframe.
%   - Quantitative claims (fill after results).
%   - Open-source call to action.

We have presented VAJRA, an open-source LLM red-teaming
framework that replaces heuristic topic-keyword scoring with
mechanically verified ground truth via canary-token exfiltration
proofs and sandboxed exploit execution, and with a decomposed
rubric judge validated against a published consensus-labeled
prompt bank. VAJRA is released at
\url{https://github.com/sahilll05/LLM-RedTeam}. All ingestion
scripts, scoring engine code, and canary-token placement
variants are fully reproducible; gated dataset access requires
the user's own Hugging Face token.

% ============================================================
\bibliographystyle{IEEEtran}
\bibliography{references}

\end{document}
